\documentclass[
  numbers=section,
  document-code={ВКРБ-09.03.04-10.19-03-24-81}
]{../fqw}

\begin{document}

\FqwEnableDocumentCodeHeader

\tableofcontents

\FqwUnnumberedSection{Введение}

Текст введения без сокращений. Иллюстрация приведена в соответствии с рисунком~\ref{fig:test}.

\section{Анализ предметной области}

Основной текст должен быть выровнен по ширине и иметь абзацный отступ.
К основным элементам относятся:
\begin{itemize}
\item первый элемент;
\item второй элемент.
\end{itemize}

\begin{figure}[H]
  \centering
  \fbox{\rule{0pt}{20mm}\rule{40mm}{0pt}}
  \caption{Проверочный рисунок}
  \label{fig:test}
\end{figure}

Параметры приведены в таблице~\ref{tab:test}.

\begin{table}[H]
  \caption{Проверочная таблица}
  \label{tab:test}
  \centering
  \begin{tabular}{|c|c|}
    \hline
    Параметр & Значение \\
    \hline
    A & 1 \\
    \hline
  \end{tabular}
\end{table}

Расчёт выполняется по формуле~\eqref{eq:test}:
\begin{equation}
  x = y + z.
  \label{eq:test}
\end{equation}
\begin{FqwWhere}
  $x$ -- значение параметра; \\
  $y$ и $z$ -- составляющие.
\end{FqwWhere}

\FqwUnnumberedSection{Заключение}

Заключительный текст.

\appendix
\FqwAppendix{Проверочное приложение}

Текст приложения.

\end{document}
